\begin{resumo}

% Rascunho factual (2026-10-22), revisado em 2026-10-29 (parecer de orientação) e em
% 2026-10-30 (análise crítica):
% - parágrafo único, 3ª pessoa, objetivo → método → resultados → conclusão (NBR 6028);
% - alegação: desidentificação facial, não anonimização (LGPD art. 12; Seção~\ref{sec:ameaca});
% - substituição híbrida: difusão só na identidade-fonte; pixels da troca de rosto (F2);
% - ponto de operação β 1,2 (não o braço mais forte); custo em emoção explícito (G6);
% - comparação pelo par (rank-1 entre rostos com face sintética, parcela coberta): o
%   rank-1 do vídeo final premia ocultar (F1).
Vídeos de sala de aula expõem a identidade de alunos e professores. Este
trabalho investiga se trocar cada rosto detectado por uma identidade sintética
reduz a reidentificação facial automática e mantém o sinal comportamental
(expressão, pose, emoção), medido no vídeo que o espectador vê. O pipeline é
híbrido e tem três estágios: detecção SCRFD-10GF com rastreamento e um pós-passe
de recall; uma identidade sintética Stable Diffusion XL por track, trocada no
rosto de cada quadro por \texttt{inswapper} com restauração GFPGAN; e composição
\emph{fail-closed}, que oculta todo rosto detectado sem substituição válida. A
contribuição principal empurra o embedding de identidade usado pela troca para
longe da identidade real estimada no track. Em três vídeos curtos (29\,s no
total), no ponto de operação escolhido, o rank-1 do FaceNet, reconhecedor não
usado na otimização, cai de \result{e3c/pooled/native/facenet.rank1} para
\result{e3c/pooled/b-1.2/facenet.rank1}, sem diferença mensurável em expressão e
pose e com queda não significativa da concordância de emoção
(\result{e3c/pooled/native/emotion_kappa} para \result{e3c/pooled/b-1.2/emotion_kappa}).
Entre os rostos que recebem face sintética, o pipeline final é reidentificado em
\result{derived/frontier/e5-ref/rank1} cobrindo \result{derived/frontier/e5-ref/share}
dos rostos, contra \result{derived/frontier/e2-blanket/rank1} do BLANKET liberado,
que cobre \result{derived/frontier/e2-blanket/share}; no vídeo final, porém,
métodos que ocultam mais rostos têm rank-1 menor, e borramento, mosaico e GANs
têm rank-1 menor que o do pipeline. Orientação por gradiente no amostrador de difusão e
otimização do código de identidade de uma rede adversarial generativa não têm
efeito mensurável. O erro de expressão fica em
\result{e3e/pooled/gate/expression} unidades de distância interocular, cerca de
quatro vezes o piso de ruído, e nenhum método preserva a emoção. O custo projetado
do laço por rosto é de \result{e5/pooled/adapt-0.08/class-hours} horas de unidade
de processamento gráfico (GPU) por hora de aula. O resultado é desidentificação
facial parcial, outro ponto do compromisso entre privacidade, cobertura e
utilidade, e não anonimização: cabelo, contorno da cabeça, roupa e voz permanecem.

\end{resumo}
